\section{Background}
\subsection{Data Visualisation \& Data Storytelling}
Sensemaking is a key process through which humans interpret and ascribe meaning to information \cite{campos2021makingsense} relying on mental models that serve as functional internal representations of external systems. These mental models preserve structural, behavioural, and data-related properties, enabling reasoning and problem-solving through simulation in working memory \cite{liu2010mental}. \citet{klein2007data} describe sensemaking as a ``\textit{deliberate effort to understand events}'' with key functions including forming explanations, predicting future states, and identifying relationships. Data visualisations, such as line charts or bar charts, support sensemaking by enabling \textit{exploratory} analysis to uncover patterns and insights \cite{Zhang2020dataquilt, knaflic2015storytelling, evergreen2020effective}. These exploratory visualisations are often minimalistic and tailored for audiences familiar with the data context or equipped with analytical skills \cite{martinezmaldonado2020layeredstorytelling, srinivasan2019augmenting}.

Existing research has recognised the potential role that narratives can have on the sensemaking process \cite{figueiras2014narrative, pirolli2005sensemaking}, particularly in aligning external visual representations with users’ mental models \cite{liu2010mental}. To address diverse levels of visualisation literacy and facilitate cognitive offloading \cite{liu2010mental,hullman2011visualization}, \textit{explanatory} visualisations have evolved, incorporating storytelling elements to communicate insights effectively \cite{figueiras2014narrative}. These author-driven “\textit{data stories}” guide audiences through structured narratives that emphasise key findings \cite{segel2010narrative, hullman2011visualization, echeverria2018exploratory}. Data storytelling (DS) combines data, visual elements, and narratives to contextualise and share insights \cite{ryan2016visual, dykes2020effective}. Principles outlined by \citet{echeverria2018exploratory} and \citet{martinezmaldonado2020layeredstorytelling} emphasise guiding audience attention through intentional visualisation, narrative choices, and alignment with specific goals. Complementing this, \citet{hullman2011visualization} argues for the use of rhetorical devices such as framing, annotations, and metonymy guide interpretation, focusing the audience’s attention on key insights to maintain a visual narrative flow \cite{mckenna2017visual}.

DS can take various formats, such as data comics \cite{bach2017emerging, bach2018designpatterns, milesi2024comics}, narrative slideshows \cite{chen2019designing}, data videos \cite{amini2015datavideos, xian2023datavideos}, and annotated charts \cite{echeverria2018exploratory}. Among these, narrative-enhanced visualisations stand out for their effectiveness. Visual features of this format often align with \citet{tufte2001visual}'s “\textit{data-ink ratio},” emphasising components that directly support insight delivery \cite{knaflic2015storytelling, zdanovic2022influence}. Design features proposed by \citet{knaflic2015storytelling} include: 
\begin{enumerate}
    \item \textit{\textbf{Highlighting important data points}} by using visual attributes like colour, shading, or shapes to direct attention to critical insights \cite{ryan2016visual,hullman2011visualization}.
    \item \textit{\textbf{Explanatory titles}} that guide the interpretation of the visualisation by directly communicating the main insight intended by the creator \cite{borkin2016beyond, echeverria2018exploratory}. 
    \item \textit{\textbf{Annotations}} that provide context or explain key takeaways, focusing viewers on significant aspects \cite{hullman2011visualization, knaflic2015storytelling}.
\end{enumerate}

Additionally, DS-enhanced visualisations often employ \textit{\textbf{decluttering}} to simplify the communication of potentially complex ideas by removing non-essential elements, thereby minimising cognitive load and clarifying the narrative \cite{hullman2011visualization, knaflic2015storytelling}.

\subsection{Data Storytelling Effectiveness and Efficiency}
%% Theoretical support
Theoretical researchers and visualisation experts, such as \citet{segel2010narrative}, \citet{gershon2001}, and \citet{Daradkeh2021} have argued that combining narratives with graphics makes visualisations more intuitive and should, in theory, improve comprehension and decision-making. An explanation for this is that DS leverages pre-attentive processing, using visual attributes like colour and shape to guide attention and facilitate quicker understanding of key data points \cite{krum2013cool,Hicks2009perceptual, knaflic2015storytelling}. However, empirical researchers have found mixed evidence regarding the effectiveness and effectiveness of DS-enhanced visualisations.  

%% Effectiveness
Regarding the effectiveness in low cognitive tasks like remembering, \citet{bateman2010} found that visual embellishments, such as colour and textual pointers, did not impact short-term recall compared to minimalist visualisations but did significantly enhance long-term recall. In contrast, \citet{zdanovic2022influence} found no significant differences in memory retention between DS-enhanced and conventional visualisations in either short-term or long-term contexts. Further evidence supporting the effectiveness of DS comes from \citet{shao2024datastorytelling}, who conducted a study to explore how DS-enhancements affect information retrieval and insight comprehension. Their findings suggest that DS visualisations improve the effectiveness of comprehension tasks, particularly those involving single insights, compared to conventional visualisations. Moreover, \citet{stokes2023striking} found that charts with heavy textual annotations were preferred by users and led to better comprehension, particularly when the annotations highlighted statistical or relational components. Similarly, \citet{kong2018} found that framing the titles of visualisations influenced users' perceptions and recall, indicating that even subtle DS elements like title framing can affect how insights are derived. It is noteworthy that these studies were limited to tasks requiring low levels of cognition like finding specific data points or understanding the main take away message of the visualisation. 

%% Efficiency
The efficiency of DS-enhanced visualisations — how quickly individuals can extract insights — has also been the subject of investigation to some extent. Small-scale eye-tracking studies \cite{de2014evaluating,Pozdniakov2023how} suggest that DS elements attract users' attention more effectively, yet they do not provide direct evidence of whether such elements facilitate quicker comprehension. In an educational context, \citet{echeverria2018exploratory} found that DS elements can reduce the time required to process information, where teachers were able to interpret data faster with DS-enhanced visualisations compared to conventional ones. However, the only large-scale empirical study on efficiency comparing conventional visualisations against those enhanced with DS elements \cite{shao2024datastorytelling} found that while DS elements improved the effectiveness of insight comprehension, they did not necessarily make the process faster.

\subsection{Data Storytelling for Supporting Tasks of Varying Complexity}
The relationship between task complexity and the potential impact that it can have on the effectiveness of visualisations has already been explored in the literature. Many of these prior studies have used Bloom's taxonomy, either partially or fully, as a framework to understand this effect. Bloom's taxonomy is a hierarchical framework that emerged from educational literature \cite{bloom1956} and it has recently been adapted and adopted in information visualisation literature \cite{burns2020evaluate, leerobbins2022learning}. It organises tasks by cognitive complexity, ranging from basic levels of thinking, such as identifying and understanding, to higher-order cognitive processes including analysing, evaluating, and creating. In the revised version of Bloom's taxonomy \cite{krathwohl2002revision}, the category previously referred to as \textit{Knowledge} was renamed \textit{Remember} to reflect the focus on retrieving relevant knowledge. In the context of information visualisation (InfoVis), however, the first level is often referred to as \textit{Identify}, as it better captures the process of recognising and selecting information from visual data \cite{shao2024datastorytelling}. Accordingly, throughout this paper, we will use the term \textit{Identify} to refer to the first level of Bloom's taxonomy.

The first level of Bloom's taxonomy was the key focus of the \citet{zdanovic2022influence} 
exploration into the effects that DS has on the ability to recall information delivered by visualisations. They found that there were no significant differences in recall between conventional and DS-enhanced visualisations.

In their investigation into the effect of DS on information retrieval and insight comprehension, \citet{shao2024datastorytelling} found that there were no significant differences in how efficiently participants were able to successfully answer multiple-choice questions based on DS-enhanced visualisations compared to conventional visualisations. Conversely, they found that the addition of storytelling elements can significantly influence how effectively individuals comprehend data insights, as measured by the number of correct responses. However, this study was limited to the first two levels of Bloom's taxonomy: \textit{Remember/Identify data points} and \textit{Understand}.

In contrast to the aforementioned studies that focus primarily on the lower levels of Bloom's taxonomy, \citet{burns2020evaluate} provided a framework adapted from the full extent of the original taxonomy to evaluate the levels to which participants understood visualisations. Although the visualisations do not explicitly incorporate the design features of DS from \citet{knaflic2015storytelling}, \citet{burns2020evaluate} highlights that the design decisions made when creating visualisations can influence a reader's ability to correctly interpret information. This motivates the exploration of the effectiveness and efficiency of visualisation with DS elements, not only at lower levels of cognitive complexity but also to understand the extent to which the design enhancements impact higher-level tasks.

Recently, the revised Bloom's taxonomy from \citet{krathwohl2002revision} has been used by \citet{adar2021communicative} and \citet{leerobbins2022learning} as part of their research into an explanatory approach called communicative visualisations. Specifically, Bloom's taxonomy inspired the structure of learning objectives, a method of framing the intention of the visualisation in order to guide the audience, as well as evaluate the "success" of the visualisation's design \cite{adar2021communicative}. However, while the effectiveness of the visualisation design is discussed using metrics such as memorability, readability, and user preference, there is minimal discussion on the relationship between the complexity of the task (as measured by Bloom's taxonomy) and the design of the visualisations. 
